\begin{frame}{한눈에 보기: 같은 특징 추출기, 다른 헤드}
  \small
  \begin{tabular}{@{}lllp{5.6cm}@{}}
    \toprule
    문제 & 출력 형태 & 특징 추출기 위에 붙는 헤드 & 대표 모델 \\
    \midrule
    분류 & 벡터 1개(K개 클래스 점수)
        & 격자 펼치기 + 선형층 1개 & VGG, ResNet \\
    객체 탐지 & (클래스, 박스 좌표 4개) $\times$ N개 사물
        & 후보 생성 + 박스별 분류기(SVM) + NMS
        & R-CNN, Fast R-CNN, YOLO \\
    의미론적 분할 & 입력과 같은 크기의 클래스 지도
        & 위치별 \textbf{공유 1$\times$1 합성곱}
          (출력 채널 = 클래스 수) & U-Net, FCN \\
    전이학습 & (위 중 아무거나)
        & \textbf{헤드만} 재학습, 또는 낮은 학습률로 전체
          미세조정 & — \\
    \bottomrule
  \end{tabular}
  \vskip1em
  \begin{block}{장 마무리}
    네 줄 전부 10.2의 ``\textbf{합성곱 스택}''이라는 \textbf{같은
    줄기에 다른 헤드}를 붙인 것 --- 헤드의 차이에 따라 답의
    \textbf{형태만} 바뀐다(벡터 $\to$ 박스 목록 $\to$ 픽셀 지도).
    CNN은 ``\textbf{이미지의 국소적 패턴은 위치와 무관하게
    재사용될 수 있다}''는 하나의 통찰을, \textbf{합성곱이라는
    연산 하나로 구현}한 것 --- 전이학습·탐지·분할·ResNet은
    서로 다른 문제를 풀지만 전부 그 재료를
    ``\textbf{어떻게 재사용·재조합·안정화할 것인가}''를
    답하는 응용.
  \end{block}
\end{frame}
